\documentclass[11pt,a4paper]{article}
\usepackage[a4paper,margin=2.25cm]{geometry}
\usepackage[T1]{fontenc}
\usepackage{lmodern}
\usepackage{graphicx,booktabs,amsmath,amssymb,xcolor,array,tabularx,float}
\usepackage{microtype}
\usepackage{fancyhdr}
\usepackage{hyperref}
\definecolor{navy}{HTML}{12314A}
\definecolor{teal}{HTML}{176B87}
\definecolor{softblue}{HTML}{EDF5F8}
\definecolor{softgreen}{HTML}{EFF7F0}
\definecolor{softpeach}{HTML}{FFF3E9}
\hypersetup{colorlinks=true,linkcolor=navy,urlcolor=teal,citecolor=teal,
  pdftitle={Estimation Theory and Regression Inference: Phase 2},
  pdfauthor={Mahdi Abazari Aghdam}}
\pagestyle{fancy}\fancyhf{}
\fancyhead[L]{\small Estimation Theory and Regression Inference}
\fancyhead[R]{\small Phase 2 Portfolio}\fancyfoot[C]{\thepage}
\setlength{\headheight}{15pt}
\setlength{\parskip}{5pt}\setlength{\parindent}{0pt}
\newcommand{\sectionpage}[2]{\clearpage\pagecolor{#1}\section{#2}}
\newcommand{\fig}[3]{\begin{figure}[H]\centering\includegraphics[width=#2\linewidth]{../figures/#1.png}\caption{#3}\label{fig:#1}\end{figure}}
% Generated by build_report.py; do not edit values by hand.
\newcommand{\LifeN}{60}
\newcommand{\LifeMean}{356.950}
\newcommand{\LifeMedian}{237.400}
\newcommand{\LifeSD}{429.883}
\newcommand{\LifeMin}{6.800}
\newcommand{\LifeMax}{2468.600}
\newcommand{\LifeCV}{1.204}
\newcommand{\ThetaMoM}{356.950}
\newcommand{\ThetaMLE}{356.950}
\newcommand{\BiasTrials}{10000}
\newcommand{\MCMean}{356.683}
\newcommand{\MCBias}{-0.267}
\newcommand{\MCVariance}{2084.362}
\newcommand{\TheoryVariance}{2123.555}
\newcommand{\CRLB}{2123.555}
\newcommand{\Efficiency}{1.019}
\newcommand{\ConsistencyTrials}{3000}
\newcommand{\AsymptoticQQ}{0.999}
\newcommand{\ExactWidth}{186.349}
\newcommand{\BootstrapWidth}{215.793}
\newcommand{\BootData}{5000}
\newcommand{\Reliability}{0.246}
\newcommand{\ThetaNumeric}{356.950}
\newcommand{\ExactCI}{[281.411, 467.760]}
\newcommand{\BootCI}{[261.169, 476.962]}
\newcommand{\ReliabilityCI}{[0.147, 0.351]}
\newcommand{\CoverageTrials}{2000}
\newcommand{\CoverageBoot}{1000}
\newcommand{\ExactCoverage}{0.9510}
\newcommand{\BootCoverage}{0.9375}
\newcommand{\VarEmpTen}{12771.634}
\newcommand{\VarTheoryTen}{12741.330}
\newcommand{\VarEmpFifty}{2503.263}
\newcommand{\VarTheoryFifty}{2548.266}
\newcommand{\VarEmpTwoHundred}{650.906}
\newcommand{\VarTheoryTwoHundred}{637.067}
\newcommand{\VarEmpThousand}{124.364}
\newcommand{\VarTheoryThousand}{127.413}
\newcommand{\VarEmpFiveThousand}{25.365}
\newcommand{\VarTheoryFiveThousand}{25.483}
\newcommand{\ProcessN}{40}
\newcommand{\Intercept}{108.602}
\newcommand{\Slope}{0.73332}
\newcommand{\TempMean}{349.675}
\newcommand{\HardnessMean}{365.025}
\newcommand{\Sxx}{545178.775}
\newcommand{\Sxy}{399788.925}
\newcommand{\Syy}{313954.355}
\newcommand{\SSE}{20782.295}
\newcommand{\ErrorVariance}{546.902}
\newcommand{\ResidualSE}{23.386}
\newcommand{\RSquared}{0.93380}
\newcommand{\InterceptNumeric}{108.602}
\newcommand{\SlopeNumeric}{0.73332}
\newcommand{\SigmaMLE}{519.557}
\newcommand{\RegDF}{38}
\newcommand{\TCrit}{2.024}
\newcommand{\InterceptSE}{11.676}
\newcommand{\SlopeSE}{0.03167}
\newcommand{\SlopeT}{23.153}
\newcommand{\SlopeP}{\ensuremath{5.24\times10^{-24}}}
\newcommand{\InterceptCI}{[84.965, 132.239]}
\newcommand{\SlopeCI}{[0.66920, 0.79744]}
\newcommand{\ShapiroW}{0.971}
\newcommand{\ShapiroP}{0.391}
\newcommand{\BPLM}{0.546}
\newcommand{\BPP}{0.460}
\newcommand{\DW}{1.593}
\newcommand{\PredFourHundred}{401.929}
\newcommand{\MeanCI}{[393.778, 410.081]}
\newcommand{\NewPartPI}{[353.890, 449.968]}
\newcommand{\TempMin}{172.000}
\newcommand{\TempMax}{553.000}
\newcommand{\RegMCTrials}{10000}
\newcommand{\RegInterceptBias}{0.146}
\newcommand{\RegSlopeBias}{\ensuremath{-4.34\times10^{-4}}}
\newcommand{\RegSlopeVarEmp}{\ensuremath{9.94\times10^{-4}}}
\newcommand{\RegSlopeVarTheory}{0.00100}
\newcommand{\RegSlopeCoverage}{0.95250}

\begin{document}
\thispagestyle{empty}
\begin{center}
\vspace*{0.6cm}
\includegraphics[width=2.8cm]{assets/sharif_logo.png}\\[0.8cm]
{\color{navy}\rule{\linewidth}{1.7pt}}\\[0.8cm]
{\Huge\bfseries Estimation Theory}\\[0.2cm]
{\Huge\bfseries \& Regression Inference}\\[0.4cm]
{\LARGE Phase 2: A From-Scratch Statistical Engine}\\[0.9cm]
{\large Reliability Engineering \(\cdot\) Process Regression\\
Simulation \(\cdot\) Confidence Intervals \(\cdot\) Prediction}\\[0.9cm]
{\color{navy}\rule{\linewidth}{1.7pt}}\\[1.3cm]
{\Large Sharif University of Technology}\\[0.3cm]
{\large Department of Industrial Engineering}\\[0.2cm]
{\large Engineering Statistics}\\[1.3cm]
{\small AUTHOR}\\[0.2cm]
{\Large\bfseries Mahdi Abazari Aghdam}\\[0.8cm]
{\large October 2026}
\end{center}
\clearpage
\pagecolor{softblue}
\section*{Abstract}
This project implements statistical estimation and regression directly from course formulas on two supplied process datasets. For 60 failure times, the method-of-moments and exponential maximum-likelihood estimates both give a mean lifetime of \ThetaMLE{} hours. A 10,000-trial simulation confirms negligible bias and variance close to the Cram\'er--Rao bound. The exact 95\% confidence interval is \ExactCI{} hours; a percentile bootstrap interval is \BootCI. In 2,000 repeated datasets, exact and bootstrap intervals cover the plug-in generating value in \ExactCoverage{} and \BootCoverage{} of trials. Estimated survival to 500 hours is \Reliability. For 40 heat-treatment observations, a directly calculated least-squares fit gives slope \Slope{} HV/\textdegree C and $R^2=\RSquared$. A Gaussian likelihood optimizer independently reproduces the fitted coefficients. At 400\textdegree C, predicted mean hardness is \PredFourHundred{} HV; the interval for one new part is much wider than the interval for the mean. All numbers and figures are regenerated by a fixed-seed Python command. The supplied files are realistic-looking assignment data, so physical generalization requires external validation.
\setcounter{tocdepth}{1}
\tableofcontents
\clearpage\nopagecolor

\section{Introduction}
The original Phase 2 assignment \cite{assignment} connects two core ideas: estimating one distribution parameter from component lifetimes, and estimating a functional relation between temperature and hardness. Both can be framed through likelihood, yet their finite-sample uncertainty is different. Part A assumes independent exponential lifetimes with mean $\theta$; Part B assumes a linear mean response with independent Gaussian errors. The project tests those assumptions where possible and distinguishes mathematical derivations, numerical answers on the supplied files, and simulation or independent-library verification. The data are provided as realistic-looking course examples; no independent provenance or production sampling design is documented.

\section{Methods and reproducibility}
Every reported point estimator and regression coefficient is implemented from direct sums in \texttt{src/estimation.py} or \texttt{src/regression.py}. Distribution quantiles, random number generation, and generic numerical optimization are allowed helpers. Fitting-library equivalents occur only in functions with a \texttt{\_verify} suffix and do not feed the reported answers. A local NumPy generator uses seed 403109076. The bias study uses \BiasTrials{} trials; the consistency study uses \ConsistencyTrials{} trials at each of five sample sizes. Dataset bootstrap intervals use \BootData{} resamples. The coverage study uses \CoverageTrials{} repeated datasets and \CoverageBoot{} bootstrap resamples per dataset. The latter count is lower to keep the full coverage experiment computationally practical. A separate regression simulation uses \RegMCTrials{} trials. All 14 figures are exported at 300 dpi. Confidence intervals use a two-sided 95\% level unless stated otherwise. The exact exponential interval has model-based coverage; a bootstrap interval has approximate coverage and depends on the empirical distribution \cite{bootstrap}.

\sectionpage{softgreen}{Part A: Exponential reliability estimation}
\subsection{A1. Exploratory description and model plausibility}
The \LifeN{} observed lifetimes have mean \LifeMean{} hours, median \LifeMedian, sample SD \LifeSD, minimum \LifeMin, and maximum \LifeMax. Their sample coefficient of variation is $s/\bar x=\LifeCV$. For an Exponential($\theta$) distribution, both the population mean and standard deviation equal $\theta$, so its population CV is exactly one. The observed CV is moderately above one, and the long upper tail in Figure~\ref{fig:01_lifetime_histogram} makes the model plausible as a working approximation, not proven. A single large lifetime can materially affect both mean and CV.
\fig{01_lifetime_histogram}{0.8}{The 60 supplied lifetimes are right-skewed; the fitted exponential density is overlaid for visual assessment.}

\subsection{A2. Method-of-moments estimate}
Equating the first sample moment to $E[X]=\theta$ gives $\hat\theta_{\mathrm{MoM}}=\bar X=\frac1n\sum_iX_i$. Direct summation yields \ThetaMoM{} hours.

\subsection{A3. Closed-form and numerical maximum likelihood}
The log likelihood is $\ell(\theta)=-n\log\theta-\sum_i x_i/\theta$. Its score vanishes at $\hat\theta_{\mathrm{MLE}}=\bar x=\ThetaMLE$ hours. The second derivative at that point is negative, as shown in Appendix~\ref{sec:derivation}. The MoM and MLE answers coincide because the sample mean is both the first empirical moment and the statistic through which the exponential scale enters the likelihood. A separate generic optimizer returns \ThetaNumeric{} hours, and a separately named SciPy fit agrees to numerical precision. Neither verification function supplies the answer.

\subsection{A4. Bias and variance, proven and simulated}
Linearity of expectation gives $E[\bar X]=\theta$; independence gives $\mathrm{Var}(\bar X)=\theta^2/n$. In \BiasTrials{} repeated 60-item samples generated at plug-in $\theta=\ThetaMLE$ hours, the mean estimate is \MCMean{} hours and empirical bias is \MCBias{} hours. Empirical variance is \MCVariance{} hours$^2$, compared with theoretical \TheoryVariance{} hours$^2$. Small differences are Monte Carlo noise, not estimator bias.
\fig{02_estimator_mc}{0.8}{Across 10,000 independent samples of size 60, estimated mean lifetimes cluster near the generating value.}

\subsection{A5. Consistency and large-sample shape}
Because the estimator is unbiased and $\mathrm{Var}(\bar X)=\theta^2/n\to0$, Chebyshev's inequality implies $\bar X\xrightarrow{p}\theta$. Table~\ref{tab:consistency} compares empirical and theoretical variance across five $n$ values, using \ConsistencyTrials{} replications per row. Figures~\ref{fig:03_consistency} and \ref{fig:04_consistency_loglog} show the expected $1/n$ decline. At $n=5000$, standardized estimates have Q--Q $R^2=\AsymptoticQQ$, consistent with the central limit theorem.
\begin{table}[H]\centering\caption{Monte Carlo variance of $\hat\theta$ approaches $\theta^2/n$; 3,000 trials per sample size.}\label{tab:consistency}
\begin{tabular}{rrr}\toprule $n$ & Empirical variance (hours$^2$) & Theoretical variance (hours$^2$)\\\midrule
10 & \VarEmpTen & \VarTheoryTen\\
50 & \VarEmpFifty & \VarTheoryFifty\\
200 & \VarEmpTwoHundred & \VarTheoryTwoHundred\\
1000 & \VarEmpThousand & \VarTheoryThousand\\
5000 & \VarEmpFiveThousand & \VarTheoryFiveThousand\\\bottomrule\end{tabular}\end{table}
\fig{03_consistency}{0.76}{Empirical estimator variance follows the theoretical $\theta^2/n$ curve.}
\fig{04_consistency_loglog}{0.76}{On log--log axes, estimator variance follows approximately a slope of $-1$.}
\fig{05_asymptotic_qq}{0.76}{At $n=5000$, standardized Monte Carlo estimates are approximately normal.}

\subsection{A6. Cram\'er--Rao bound and efficiency}
One exponential observation has Fisher information $1/\theta^2$, so $n$ independent observations have information $n/\theta^2$. The variance lower bound for an unbiased regular estimator is therefore $\theta^2/n=\CRLB$ hours$^2$. The sample mean has exactly this theoretical variance and is efficient under the assumed model. The numerical ratio $\mathrm{CRLB}/\widehat{\mathrm{Var}}_{MC}=\Efficiency$ is slightly above one because the empirical variance is itself noisy; it does not contradict the bound.
\fig{06_efficiency}{0.68}{The 10,000-trial empirical variance is close to the Cram\'er--Rao bound.}

\subsection{A7. Exact versus bootstrap confidence intervals}
The pivot $2n\bar X/\theta\sim\chi^2_{2n}$ gives an exact 95\% interval of \ExactCI{} hours (width \ExactWidth). A \BootData-resample nonparametric percentile interval is \BootCI{} hours (width \BootstrapWidth). The bootstrap is wider in this observed sample. It can be useful when no pivotal formula is available, but it does not inherit exact coverage from the exponential model.
\fig{07_bootstrap_lifetime}{0.76}{The 5,000-resample distribution of mean lifetime shows the percentile interval around the observed estimate.}
In \CoverageTrials{} repeated 60-item datasets generated at the plug-in value, the exact interval covered the generating parameter in \ExactCoverage{} of trials and the bootstrap interval in \BootCoverage. The latter used \CoverageBoot{} resamples per dataset. Figure~\ref{fig:08_coverage} compares both to the nominal .95 target. This experiment checks the procedures under an assumed exponential generator; it does not establish that the real production mechanism is exponential.
\fig{08_coverage}{0.7}{Exact coverage is near 95\%; percentile bootstrap coverage is modestly lower across 2,000 model-generated samples.}

\subsection{A8. Reliability at 500 hours}
The exponential survival function is
\begin{equation*}R(t)=\exp(-t/\theta),\qquad
\hat R(500)=\exp(-500/\hat\theta)=\Reliability.
\end{equation*}
The estimate follows from MLE invariance. Transforming each of the \BootData{} bootstrap mean estimates gives a 95\% interval \ReliabilityCI. This is a model-based survival probability at 500 hours, not a directly observed proportion of components in the supplied sample.
\fig{09_reliability}{0.75}{Estimated survival probability declines with operating time; the 500-hour estimate is highlighted.}

\sectionpage{softpeach}{Part B: Temperature and hardness regression}
\subsection{B1. Scatterplot and linearity}
The file contains \ProcessN{} pairs spanning \TempMin--\TempMax\textdegree C. Figure~\ref{fig:10_process_fit} shows a strong approximately linear increase in hardness over this observed range. Curvature or fanning in a residual plot would weaken that description. The file documents no randomization of treatment temperature, so the fitted association should not be called a causal effect without further process-design evidence.
\fig{10_process_fit}{0.82}{A line fitted from direct sums describes the temperature--hardness trend across 40 observations.}

\subsection{B2. Least-squares fit from direct sums}
The course formulas give $S_{xx}=\sum(x_i-\bar x)^2=\Sxx$ and $S_{xy}=\sum(x_i-\bar x)(y_i-\bar y)=\Sxy$. Thus $b=S_{xy}/S_{xx}=\Slope$ HV/\textdegree C and $a=\bar y-b\bar x=\Intercept$ HV. The fitted relation is $\hat y=\Intercept+\Slope x$. An independent SciPy linear fit agrees to numerical precision; it is only a verification.

\subsection{B3. Equivalence of least squares and Gaussian MLE}
With Gaussian errors, $\ell(A,B,\sigma^2)=-\frac n2\log(2\pi\sigma^2)-SSE(A,B)/(2\sigma^2)$. At any positive variance, maximizing over $A,B$ is equivalent to minimizing SSE. Differentiation yields the same two normal equations as least squares (Appendix~\ref{sec:derivation}). A generic numerical likelihood optimizer returns intercept \InterceptNumeric{} and slope \SlopeNumeric, agreeing with the direct fit to $10^{-5}$ or better. The coefficient equivalence does not mean all variance estimates coincide: the Gaussian MLE uses $SSE/n$.

\subsection{B4. Error variance}
The direct residual sum of squares is $SSE=\SSE$ HV$^2$. Dividing by $n-2=\RegDF$ gives unbiased $s^2=\ErrorVariance$ HV$^2$, and residual standard error $s=\ResidualSE$ HV. This is a typical vertical deviation of observed hardness from the fitted line. The Gaussian variance MLE is \SigmaMLE{} HV$^2$, smaller because its divisor is $n$.

\subsection{B5. Intervals and slope hypothesis test}
Using $t_{\RegDF}$ with critical value \TCrit, the 95\% intercept interval is \InterceptCI{} HV and the 95\% slope interval is \SlopeCI{} HV/\textdegree C. The test $H_0:B=0$ against $H_1:B\ne0$ gives $t=\SlopeT$, $df=\RegDF$, and $p=\SlopeP$. The data therefore support a nonzero linear association at $\alpha=.05$. Since 0\textdegree C is outside the observed range, the intercept is a model parameter rather than an operating-point hardness prediction.
In \RegMCTrials{} new 40-part datasets generated from the fitted Gaussian line, direct-sum slope estimates show bias \RegSlopeBias{} HV/\textdegree C. Their empirical variance is \RegSlopeVarEmp, compared with theoretical $s^2/S_{xx}=\RegSlopeVarTheory$. The 95\% slope interval covered the generating slope in \RegSlopeCoverage{} of trials. These checks support the formula implementation under the plug-in model, without proving that the factory errors are Gaussian.
\fig{14_regression_mc}{0.76}{Across 10,000 plug-in Gaussian samples, direct-sum slopes cluster around the generating coefficient.}

\subsection{B6. Coefficient of determination}
With $S_{yy}=\Syy$ HV$^2$, $R^2=1-SSE/S_{yy}=\RSquared$. Thus the fitted line accounts for about 93.4\% of the hardness variation \emph{within this sample}. High $R^2$ alone does not validate the Gaussian or constant-variance assumptions.

\subsection{B7. Residual diagnostics}
Figure~\ref{fig:11_residuals_fitted} plots residuals against fitted hardness; Figure~\ref{fig:12_residuals_qq} checks normal shape. Shapiro--Wilk yields $W=\ShapiroW$, $p=\ShapiroP$; no clear normality departure is detected at .05. The Breusch--Pagan auxiliary regression gives $nR^2=\BPLM$, $p=\BPP$; no clear linear variance trend with temperature is detected. These non-rejections are not proofs of the assumptions. Durbin--Watson is \DW{} in the supplied row order. Since rows are temperature sorted rather than time ordered, that statistic has no valid temporal-autocorrelation interpretation here.
\fig{11_residuals_fitted}{0.77}{Residuals versus fitted hardness provide a visual check for curvature and changing spread.}
\fig{12_residuals_qq}{0.77}{The residual normal Q--Q plot shows no strong departure from a straight-line pattern.}

\subsection{B8. Mean response and one-part prediction at 400\textdegree C}
The fitted mean at $x_0=400$\textdegree C is $\hat y_0=\PredFourHundred$ HV. The 95\% confidence interval for the \emph{process mean} is \MeanCI{} HV. The 95\% prediction interval for the hardness of \emph{one new part} is \NewPartPI{} HV. The latter is wider because it includes individual residual variation as well as uncertainty in the fitted mean. A production engineer estimating average process hardness should use the mean interval; an engineer assessing one upcoming part should use the prediction interval. This is interpolation because 400\textdegree C falls within the observed temperature range.
\fig{13_prediction_intervals}{0.73}{At 400\textdegree C, the interval for a new part is wider than the mean-response interval.}

\sectionpage{softblue}{Cross-case discussion and conclusions}
Part A's MoM--MLE equality and Part B's LSE--MLE equality are two instances of optimizing a parametric model through sufficient summary quantities. In Part A the sample mean determines the exponential likelihood's optimum; in Part B the residual sum of squares determines the Gaussian coefficient optimum. Analytical statements and Monte Carlo checks play different roles: proof establishes the property under assumptions, while simulation exposes implementation mistakes and finite-trial variability. The efficiency ratio slightly above one illustrates that a finite Monte Carlo variance estimate is not a mathematical variance bound. The exact interval's observed coverage near .95 supports its implementation under the assumed generator; the bootstrap's lower coverage shows why an approximate interval should be checked in skewed samples. Neither check validates the physical model without external production evidence.

\begin{itemize}
\item Mean component lifetime is estimated at \ThetaMLE{} hours by both MoM and MLE; the observed CV of \LifeCV{} makes the exponential assumption plausible but uncertain.
\item The sample mean is unbiased and theoretically reaches the Cram\'er--Rao bound under the exponential model; 10,000 simulations agree within Monte Carlo error.
\item Exact and bootstrap 95\% intervals differ in width and empirical coverage; the estimated 500-hour survival probability is \Reliability.
\item Hardness rises by about \Slope{} HV per 1\textdegree C in the fitted range; the slope interval excludes zero.
\item A new-part prediction at 400\textdegree C is materially less precise than the estimated process mean at that temperature.
\end{itemize}

\section{Appendix: mathematical derivations}\label{sec:derivation}
\subsection*{Exponential MoM, MLE, information, and exact interval}
For $f(x\mid\theta)=\theta^{-1}e^{-x/\theta}$, $x>0$, $E[X]=\theta$ and $\mathrm{Var}(X)=\theta^2$. Thus the first-moment equation $\bar X=\theta$ gives $\hat\theta_{MoM}=\bar X$. For $n$ independent observations,
\begin{align*}
L(\theta)&=\theta^{-n}\exp\!\left(-\frac{\sum_i x_i}{\theta}\right),\\
\ell(\theta)&=-n\log\theta-\frac{\sum_i x_i}{\theta},\\
\ell'(\theta)&=-\frac n\theta+\frac{\sum_i x_i}{\theta^2}=0
\quad\Longrightarrow\quad\hat\theta_{MLE}=\frac1n\sum_i x_i=\bar x,\\
\ell''(\hat\theta)&=\frac n{\hat\theta^2}-\frac{2\sum_i x_i}{\hat\theta^3}
=-\frac n{\hat\theta^2}<0.
\end{align*}
Unbiasedness follows from $E[\bar X]=n^{-1}\sum_i E[X_i]=\theta$; independence gives $\mathrm{Var}(\bar X)=n^{-2}\sum_i\theta^2=\theta^2/n$. For one observation, $-E[\ell''_1(\theta)]=1/\theta^2$; information adds over independent observations, so $I_n(\theta)=n/\theta^2$ and the Cram\'er--Rao lower bound is $1/I_n=\theta^2/n$. The exact pivot $2\sum_i X_i/\theta=2n\bar X/\theta\sim\chi^2_{2n}$ implies
\begin{equation*}
P\!\left(\frac{2n\bar X}{\chi^2_{2n,1-\alpha/2}}
\le\theta\le
\frac{2n\bar X}{\chi^2_{2n,\alpha/2}}\right)=1-\alpha.
\end{equation*}
Since the survival function is $g(\theta)=e^{-t/\theta}$, MLE invariance gives $\hat R(t)=g(\hat\theta_{MLE})$.

\subsection*{Regression normal equations and Gaussian likelihood}
Let $SSE(A,B)=\sum_i(y_i-A-Bx_i)^2$. Differentiating gives
\begin{align*}
\frac{\partial SSE}{\partial A}&=-2\sum_i(y_i-A-Bx_i)=0,\\
\frac{\partial SSE}{\partial B}&=-2\sum_i x_i(y_i-A-Bx_i)=0.
\end{align*}
The first equation gives $A=\bar y-B\bar x$. Substitution into the second gives $S_{xy}-B S_{xx}=0$, so $b=S_{xy}/S_{xx}$ and $a=\bar y-b\bar x$. Under independent $N(0,\sigma^2)$ errors,
\begin{equation*}
\ell(A,B,\sigma^2)=-\frac n2\log(2\pi\sigma^2)-\frac{SSE(A,B)}{2\sigma^2}.
\end{equation*}
For any $\sigma^2>0$, the coefficient score equations are $\partial SSE/\partial A=0$ and $\partial SSE/\partial B=0$, proving coefficient MLE equals LSE. Maximizing over variance gives $\hat\sigma^2_{MLE}=SSE/n$; correcting degrees of freedom for an unbiased estimate gives $s^2=SSE/(n-2)$.

\subsection*{Coefficient and response uncertainty}
With $s^2=SSE/(n-2)$, the slope and intercept standard errors are $s_b=s/\sqrt{S_{xx}}$ and $s_a=s\sqrt{1/n+\bar x^2/S_{xx}}$. Their 95\% intervals are estimate $\pm t_{.975,n-2}$ times the respective standard error. Testing $H_0:B=0$ uses $t=b/s_b$. At $x_0$, mean-response standard error is $s\sqrt{h_0}$ with $h_0=1/n+(x_0-\bar x)^2/S_{xx}$. One future observation has standard error $s\sqrt{1+h_0}$, yielding the wider prediction interval.

\section{Software and data notes}
The analysis was verified with Python 3.13, NumPy 2.2.6, pandas 2.2.3, SciPy 1.17.1, Matplotlib 3.10.8, Seaborn 0.13.2, and nbformat 5.10.4. Exact dependency pins are in \texttt{requirements.txt}. The original assignment supplied both CSV files and specified the model assumptions \cite{assignment}. The report and notebook preserve those assumptions explicitly rather than representing the files as independently verified factory measurements.

\begin{thebibliography}{9}
\bibitem{assignment} Sharif University of Technology, \emph{Engineering Statistics: Phase 2 Project, Parameter Estimation and Simple Linear Regression}, supplied course assignment.
\bibitem{bootstrap} B. Efron and R. Tibshirani, \emph{An Introduction to the Bootstrap}, Chapman \& Hall/CRC, 1993.
\bibitem{montgomery} D. C. Montgomery and G. C. Runger, \emph{Applied Statistics and Probability for Engineers}, Wiley.
\end{thebibliography}
\end{document}
